\chapter{HRBR Blocks: Template + Slots (the Compiled Fast Path)}

Chapters 12 and 13 gave us two halves of a reactive UI runtime:

\begin{itemize}
  \item \textbf{Signals} decide \emph{what} values changed.
  \item \textbf{Scheduler lanes} decide \emph{when} to apply work.
\end{itemize}

Blocks are the third half: they decide \emph{how} to touch the DOM.

A \textbf{block} is the "fast DOM" path: it describes a stable DOM template and a set of precise write locations (slots). Updates become \emph{O(k)} slot writes instead of building a VDOM tree and diffing children.

This chapter is grounded in:

\begin{itemize}
  \item Implementation: \texttt{runtime/block.ts}
  \item Specs: \texttt{runtime/\_\_tests\_\_/block.test.ts}
  \item Slot matrix: \texttt{runtime/\_\_tests\_\_/block-slot-matrix.test.ts}
  \item Input controls: \texttt{runtime/\_\_tests\_\_/block-input-controls.test.ts}
  \item Disposal correctness: \texttt{runtime/\_\_tests\_\_/block-disposal.test.ts}
\end{itemize}

\section{A tiny contract (what a block guarantees)}

Treat this as the “promise” the runtime makes to compiled code:

\begin{enumerate}
  \item \textbf{Single stable root:} \texttt{templateHTML} must produce exactly one root element, which is cloned and appended to the host.
  \item \textbf{Exact addressing:} each slot points to a node using a \texttt{path: number[]} walk over \texttt{childNodes} indices.
  \item \textbf{Specialized patching:} slot kinds are patched by small, explicit routines (text/attr/prop/class/style/event).
  \item \textbf{No needless DOM writes:} redundant values are skipped via \texttt{Object.is}.
  \item \textbf{Safe lifecycle:} \texttt{destroy()} is idempotent, removes listeners, and after destroy further \texttt{update()} calls are ignored.
\end{enumerate}

\section{Block vocabulary}

\subsection{TemplateHTML}

A block definition includes a \texttt{templateHTML: string}.
This must produce exactly one root element.

Implementation detail: \texttt{parseTemplate(templateHTML)} uses a \texttt{<template>} element and asserts:

\begin{quote}
\texttt{templateHTML must have a single root element}
\end{quote}

The implementation is intentionally strict:

\begin{itemize}
  \item it does \texttt{templateHTML.trim()} before parsing
  \item it reads \texttt{tpl.content.firstElementChild}
  \item it asserts non-null and returns the root element
\end{itemize}

This makes “multiple root nodes” a loud error, which is what you want for compiled templates.

\subsection{Slot}

A \textbf{slot} is a metadata record describing \emph{where} and \emph{how} to write a dynamic value into a template instance.

Slots are keyed by string (slot key) and come in these kinds:

\begin{itemize}
  \item \texttt{text}: write to a Text node
  \item \texttt{attr}: set/remove attributes
  \item \texttt{prop}: set DOM properties
  \item \texttt{class}: write \texttt{class} attribute
  \item \texttt{style}: write inline styles
  \item \texttt{event}: bind events once and swap handler references
\end{itemize}

Each slot includes a \texttt{path: number[]} which is a \textbf{childNodes index walk} from the template root.

Two practical consequences:

\begin{itemize}
  \item \texttt{childNodes} counts text nodes and comments, so whitespace in templates matters.
  \item Slot paths are best generated by a compiler that controls the template serialization.
\end{itemize}

\section{Why paths (and not querySelector)?}

Paths are the “compiled fast path” address:

\begin{itemize}
  \item they’re \emph{exact} (no ambiguities)
  \item they’re cheap (depth walk, no subtree scan)
  \item they enable prefix caching across multiple slots
\end{itemize}

\section{The public API (surface area)}

The core exports in \texttt{runtime/block.ts} are:

\begin{itemize}
  \item \texttt{defineBlock(def)}: identity helper for typing and readability
  \item \texttt{mountBlock(def, host, initialValues?, options?) -> MountedBlock}
  \item \texttt{mountCompiledBlock(def, host, compiledSlots, options?) -> MountedBlock \& \{ dispose() \}}
  \item \texttt{mountNestedBlock(parent, slotKey, def, initialValues?)}
  \item \texttt{mountNestedFallback(parent, slotKey, mount)}
  \item \texttt{resolvePath(root, path)} and \texttt{resolvePathCached(root, path, slotKey, dev?)}
\end{itemize}

Reading guide:

\begin{itemize}
  \item \texttt{mountBlock} is the non-reactive primitive (mount + \texttt{update}).
  \item \texttt{mountCompiledBlock} wires signals into \texttt{update} with effects and the HRBR scheduler.
  \item \texttt{mountNestedBlock}/\texttt{mountNestedFallback} are the structural escape hatches.
\end{itemize}

\section{defineBlock(): the type boundary}

	exttt{defineBlock(def)} doesn’t transform input. It exists so "compiled" code can be explicit and type-safe.

In \texttt{runtime/block.ts}, it is literally an identity function:

\begin{quote}
	exttt{export function defineBlock(def: BlockDef): BlockDef \{ return def \}}
\end{quote}

That sounds trivial, but it creates a clean “compiler boundary” and a stable type surface.

Tests use it in every case:

\begin{verbatim}
const block = defineBlock({ templateHTML: ..., slots: ... })
\end{verbatim}

\section{mountBlock(): instantiate template + resolve slot nodes}

\texttt{mountBlock()} does the non-reactive mount:

\begin{enumerate}
  \item parse template HTML into a root element
  \item clone it to create an instance
  \item append instance root to the host
  \item resolve each slot \texttt{path} into a concrete DOM \texttt{Node}
  \item apply \texttt{initialValues} via \texttt{block.update(initialValues)}
\end{enumerate}

The key optimization is “resolve once, write many”: paths are resolved at mount time and stored in \texttt{slotNodes}. Updates never traverse the DOM tree; they dereference and write.

The returned \texttt{MountedBlock} includes:

\begin{itemize}
  \item \texttt{root}: the mounted root element
  \item \texttt{slotNodes}: mapping from slot key to resolved node reference
  \item \texttt{update(values)}: apply slot writes
  \item \texttt{destroy()} / \texttt{dispose()} lifecycle
\end{itemize}

\subsection{Path resolution and caching}

\texttt{resolvePath(root, path)} is the simple implementation.
But \texttt{mountBlock} uses \texttt{resolvePathCached(instanceRoot, path, slotKey, dev)}.

The cache is a \texttt{Map<string, Node>} stored on the instance root as \texttt{\_\_hrbrPathCache}.
Prefixes are stored as dot-joined strings like \texttt{"1.0"}.

The caching algorithm is incremental: as it walks the path, it caches each prefix result (\texttt{"1"}, then \texttt{"1.0"}, etc.). Multiple slots that share a prefix benefit immediately.

Spec: \texttt{mountBlock caches repeated path prefixes across slots}.

\subsection{Dev diagnostics}

When \texttt{options.dev === true}, the assertions include the slot key in the error.

Spec: \texttt{dev diagnostics: mountBlock includes slot key in invalid path errors (dev mode)}.

\section{update(values): slot patching rules}

The update loop iterates over provided keys only. If a key doesn’t exist in \texttt{def.slots}, it is ignored.

Two important performance rules:

\begin{itemize}
  \item The block stores \texttt{prevValues[key]} and skips redundant writes using \texttt{Object.is}.
  \item It emits a devtools event \texttt{slotWrite} for each actual write.
\end{itemize}

The use of \texttt{Object.is} is deliberate: it treats \texttt{NaN} as equal to itself and provides a cheap “same reference” check for objects (like style objects or class arrays).

Spec: \texttt{skips redundant writes for unchanged values (text + attr)}.

\subsection{text slots}

Text slots must resolve to a real Text node. The runtime asserts:

\begin{quote}
\texttt{slot 'x' expected a Text node}
\end{quote}

Writes use \texttt{Text.nodeValue} and normalize nullish values to the empty string.

Spec: \texttt{mounts a template and updates text slots by path}.

\subsection{attr slots (including boolean attributes and SVG xlink)}

Attribute writes route through \texttt{setAttr(el, name, next)}:

\begin{itemize}
  \item \texttt{null | false} removes the attribute
  \item boolean-ish attributes are present/absent (set to empty string)
  \item SVG \texttt{xlink:*} attributes use \texttt{setAttributeNS}
\end{itemize}

Implementation notes:

\begin{itemize}
  \item SVG detection uses \texttt{namespaceURI} (not \texttt{instanceof SVGElement}) for broader runtime compatibility.
  \item Boolean attribute names are normalized (\texttt{readOnly} maps to \texttt{readonly}).
\end{itemize}

Specs:

\begin{itemize}
  \item \texttt{updates attribute slots by path}
  \item \texttt{treats boolean attributes as present/absent (attr slots)}
  \item \texttt{sets xlink:* attributes with the xlink namespace on SVG elements}
\end{itemize}

\subsection{prop slots (input correctness)}

Property slots route through \texttt{setProp(el, name, next)}.

Special-case correctness for inputs:

\begin{itemize}
  \item \texttt{value}: written as a string property, nullish clears to \texttt{""}
  \item \texttt{checked}: written as boolean
\end{itemize}

This matches standard controlled-input behavior: writing to properties updates the live control state, while attributes can drift.

Specs:

\begin{itemize}
  \item \texttt{updates property slots by path}
  \item \texttt{updates input value/checked via property semantics for prop slots}
  \item input control suite in \texttt{block-input-controls.test.ts}
\end{itemize}

\subsection{class slots}

Class slots write the \texttt{class} attribute:

\begin{itemize}
  \item \texttt{null | false} removes the attribute
  \item arrays are joined with spaces after filtering falsy values
  \item otherwise coerced to string
\end{itemize}

\subsection{style slots}

Style slots accept:

\begin{itemize}
  \item \texttt{null | false}: remove whole style attribute
  \item string: set \texttt{style} attribute directly
  \item object: per-key \texttt{style.setProperty} / \texttt{removeProperty}
\end{itemize}

This small matrix is validated by \texttt{block-slot-matrix.test.ts}.

\subsection{event slots: bind once, swap handler}

Event slots are designed to avoid listener leaks and avoid churn:

\begin{itemize}
  \item the first time a slot is set, the runtime binds a stable listener function
  \item the listener reads \texttt{rec.current} to find the latest handler
  \item updates replace \texttt{rec.current} without re-binding
  \item destroying the block removes the listener
\end{itemize}

Specs:

\begin{itemize}
  \item \texttt{supports event slots: binds once and updates handler without leaking listeners}
  \item \texttt{destroy removes event listeners (no further handler calls)}
  \item \texttt{skips redundant event listener churn when handler identity is unchanged}
\end{itemize}

\subsubsection{Matching Relax VDOM semantics: hostComponent binding}

\texttt{mountBlock(..., \{ hostComponent \})} optionally calls handlers with \texttt{this === hostComponent}, matching Chapter 7’s event binding behavior in the VDOM runtime.

\section{Lifecycle: destroy vs dispose, and idempotency}

\texttt{mountBlock} is not reactive by itself, so \texttt{dispose()} is just an alias for \texttt{destroy()}.

\texttt{destroy()} is idempotent:

\begin{itemize}
  \item it guards with a \texttt{destroyed} flag
  \item it destroys composed children
  \item it removes event listeners
  \item it removes the block root from the DOM
\end{itemize}

After \texttt{destroy()}, \texttt{update()} becomes a no-op (guarded by the same \texttt{destroyed} flag). That prevents stale scheduled work from mutating detached DOM.

Specs:

\begin{itemize}
  \item \texttt{destroy is idempotent even with event slots and children}
  \item \texttt{lifecycle: destroy() is idempotent and prevents further updates}
\end{itemize}

\section{mountCompiledBlock(): wiring signals to slot writes}

\texttt{mountCompiledBlock} is the bridge between the reactive graph and the block DOM writer.

Input: \texttt{CompiledSlot[]} where each entry is:

\begin{verbatim}
{ key: string, read: () => unknown }
\end{verbatim}

\textbf{Important:} \texttt{read()} can access signals; that’s how dependency tracking happens.

Mechanism:

\begin{itemize}
  \item compute initial values by running each \texttt{read()}
  \item mount a non-reactive block with those initial values
  \item create one \texttt{createEffect()} per compiled slot
  \item coalesce multiple invalidations into a single scheduled \texttt{flushPending}
  \item execute \texttt{block.update(values)} inside \texttt{batch()} to keep effects consistent
\end{itemize}

The coalescing logic (\texttt{pending} + \texttt{scheduled} flag) is the performance heart: many effects can invalidate in one turn, but only one scheduled task drains \texttt{pending} and commits to the DOM.

Scheduling goes through the deterministic lane scheduler from Chapter 13:

\begin{quote}
	exttt{scheduler.schedule(lane, flushPending)}
\end{quote}

Spec: \texttt{changing a signal updates only the intended DOM node (mountCompiledBlock)}.

That test also checks identity stability:

\begin{itemize}
  \item the static DOM elements are not recreated
  \item text node identity remains stable; only \texttt{nodeValue} changes
\end{itemize}

Spec: \texttt{mounting does not recreate static DOM during updates}.

\subsection{Disposal stops reactive updates}

\texttt{mountCompiledBlock} returns a \texttt{dispose()} method that:

\begin{itemize}
  \item disposes all effects
  \item destroys the block DOM
  \item is idempotent
\end{itemize}

Specs:

\begin{itemize}
  \item \texttt{mountCompiledBlock.dispose stops reactive updates and removes DOM}
  \item \texttt{lifecycle: mountCompiledBlock.dispose() stops reactive updates and is idempotent}
\end{itemize}

\section{Composition: nested blocks and fallback regions}

Blocks are stable and fast, but real UIs still need dynamic structure.
\texttt{runtime/block.ts} provides two small composition helpers:

\begin{itemize}
  \item \texttt{mountNestedBlock(parent, slotKey, def)}: mount a child block into an Element slot
  \item \texttt{mountNestedFallback(parent, slotKey, mount)}: mount a structural fallback region into an Element slot
\end{itemize}

Both helpers:

\begin{itemize}
  \item clear the host element’s existing children
  \item mount into it
  \item register a child destructor so destroying the parent destroys the child
\end{itemize}

Specs in \texttt{runtime/\_\_tests\_\_/block.test.ts}:

\begin{itemize}
  \item \texttt{block composition: can mount a nested block ... and destroys it with parent}
  \item \texttt{block composition: can mount a nested fallback region ... and disposes it with parent}
\end{itemize}

This is HRBR's interop story:

\begin{quote}
Compiled stable subtrees use blocks; dynamic subtrees route to fallback reconciliation.
\end{quote}

Implementation detail: composition is tracked with a small internal child registry.
	exttt{mountBlock} exposes an internal \texttt{\_\_hrbrTrackChild} function, and the composition helpers use it to register child destroyers so teardown order is correct.

\section{Walking the tests (the executable spec)}

\subsection{\texttt{runtime/\_\_tests\_\_/block.test.ts}}

This test file locks down the baseline behavior of each slot type and of the compiled/reactive mount.

\begin{itemize}
  \item \textbf{text slots:} the slot must resolve to a real Text node; updates rewrite \texttt{nodeValue}; destroy removes DOM.
  \item \textbf{attr slots:} \texttt{null} removes the attribute.
  \item \textbf{prop slots:} input \texttt{value} is updated via property semantics.
  \item \textbf{class/style:} \texttt{null} removes attrs, arrays join for class, style objects set properties.
  \item \textbf{event slots:} click once, swap handler, click again; after \texttt{destroy()} clicks do nothing.
  \item \textbf{boolean attributes:} present/absent behavior is asserted with \texttt{disabled}.
  \item \textbf{SVG xlink:} uses \texttt{getAttributeNS} to validate namespace behavior.
  \item \textbf{path caching:} asserts \texttt{\_\_hrbrPathCache} contains prefix keys (e.g. \texttt{"0"}, \texttt{"0.0"}).
  \item \textbf{dev diagnostics:} invalid path errors include the slot key in dev mode.
  \item \textbf{mountCompiledBlock:} changing one signal updates only its text node, and node identity is stable across updates.
\end{itemize}

\subsection{\texttt{runtime/\_\_tests\_\_/block-slot-matrix.test.ts}}

This is a matrix-style smoke test that mounts one template with \textbf{all slot kinds} and asserts both initial values and updated values, including event swapping.

\subsection{\texttt{runtime/\_\_tests\_\_/block-input-controls.test.ts}}

This suite documents the intended controlled-input semantics for \texttt{value} and \texttt{checked} when patched as properties.
It also includes a note about radio group behavior in jsdom.

\subsection{\texttt{runtime/\_\_tests\_\_/block-disposal.test.ts}}

This suite locks down lifecycle safety:

\begin{itemize}
  \item \texttt{destroy()} removes event listeners (clicking a detached element must not call handlers).
  \item \texttt{destroy()} is idempotent.
  \item \texttt{mountCompiledBlock.dispose()} stops reactive updates and removes DOM permanently.
\end{itemize}

\section{Online resources}

\begin{itemize}
  \item \href{https://developer.mozilla.org/en-US/docs/Web/HTML/Element/template}{MDN: \texttt{<template>}} (template parsing and cloning)
  \item \href{https://developer.mozilla.org/en-US/docs/Web/API/Node/childNodes}{MDN: \texttt{Node.childNodes}} (why paths include text/comment nodes)
  \item \href{https://developer.mozilla.org/en-US/docs/Web/API/EventTarget/addEventListener}{MDN: \texttt{addEventListener}} (listener identity + removal)
  \item \href{https://developer.mozilla.org/en-US/docs/Web/API/HTMLInputElement}{MDN: \texttt{HTMLInputElement}} (property semantics for \texttt{value} and \texttt{checked})
  \item \href{https://developer.mozilla.org/en-US/docs/Web/SVG/Attribute/xlink:href}{MDN: \texttt{xlink:href}} (SVG namespaced attribute background)
  \item \href{https://svelte.dev/blog/virtual-dom-is-pure-overhead}{Svelte: “Virtual DOM is pure overhead”} (compiler-driven DOM updates intuition)
\end{itemize}

\section{Takeaways}

\begin{itemize}
  \item Blocks are "template + slots": a stable DOM instance plus precise write locations.
  \item Slot patching is specialized per kind, with correctness for input controls and SVG namespaces.
  \item Updates skip redundant writes via \texttt{Object.is}.
  \item Event slots bind once and swap handler references; destroy removes listeners.
  \item \texttt{mountCompiledBlock} wires signals to slots and coalesces updates into scheduled flushes.
  \item Composition helpers let blocks nest blocks and fallback regions safely.
\end{itemize}
